\documentclass{article}
\usepackage[margin=1in]{geometry}
\usepackage{graphicx}
\usepackage{amsmath,amssymb,amsthm,mathtools}
\usepackage{bm}
\usepackage{booktabs}
\usepackage{enumitem}
\usepackage{hyperref}
\usepackage{listings,xcolor}

\newcommand{\R}{\mathbb{R}}
\newcommand{\C}{\mathbb{C}}
\newcommand{\E}{\mathbb{E}}
\newcommand{\calO}{\mathcal{O}}
\DeclareMathOperator{\diag}{diag}
\DeclareMathOperator{\tr}{tr}
\DeclareMathOperator{\Tr}{Tr}
\DeclareMathOperator{\Var}{Var}

\definecolor{codebg}{rgb}{0.97,0.97,0.97}
\definecolor{codekw}{rgb}{0.00,0.30,0.60}
\definecolor{codestr}{rgb}{0.58,0.00,0.00}
\definecolor{codecom}{rgb}{0.30,0.50,0.30}
\lstdefinestyle{py}{
  language=Python, backgroundcolor=\color{codebg},
  basicstyle=\ttfamily\footnotesize, keywordstyle=\color{codekw}\bfseries,
  stringstyle=\color{codestr}, commentstyle=\color{codecom}\itshape,
  showstringspaces=false, breaklines=true, columns=fullflexible,
  frame=single, framerule=0.4pt, xleftmargin=6pt, xrightmargin=6pt,
  aboveskip=4pt, belowskip=4pt,
}
\lstset{style=py}

\title{Pólya--Gamma Augmentation in GP Classification\\using Equispaced Fourier Features}
\author{Cole Citrenbaum}
\date{}

\begin{document}

\maketitle
\tableofcontents

\section{Background}
\label{sec:background}

\subsection{GP setup}
Assume $(x_i,y_i)\in\R^d\times\{0,1\}$ are i.i.d.\ observations drawn from
\[
  y_i \mid f, x_i \;\sim\; \mathrm{Bern}\!\bigl(\sigma(f(x_i))\bigr),
  \qquad
  f \;\sim\; GP(0, K_\theta).
\]
Our goal is to learn the kernel hyperparameters $\theta$ and then compute predictive means and variances. (The negative-binomial likelihood is handled identically up to likelihood-level scalars; see Section~\ref{sec:nb}.) Throughout we use the equispaced Fourier feature approximation
\[
  K \;\approx\; \Phi\Phi^{*} \;=\; F\,D^{2}\,F^{*},
  \qquad \Phi \;=\; FD,
\]
where
\begin{itemize}[leftmargin=1.4em,itemsep=1pt]
  \item $F\in\C^{n\times m}$ is a non-uniform DFT (NUFFT type-2), $[F]_{ij}=e^{i x_i\cdot\xi_j}$, evaluating $m$ Fourier modes $\xi_j$ at the $n$ data locations;
  \item $F^{*}$ is its adjoint (NUFFT type-1);
  \item $D\in\R^{m}_{\ge 0}$ is diagonal with $D_{jj}=\sqrt{h^{d}\hat k(\xi_j;\theta)}$ (quadrature-weighted square-root spectral density);
  \item $D^{2}_j=h^{d}\hat k(\xi_j;\theta)$ and $\partial_\theta D^{2}\in\R^{m\times p}$ is its gradient in the $p$ kernel hyperparameters.
\end{itemize}
Both $F$ and $F^{*}$ act in $\calO(n+m\log m)$ and are never formed as matrices.

\subsection{Pólya--Gamma augmentation}
Introduce auxiliary variables $\omega_i\sim PG(1,0)$ and use the two PG identities
\[
  \frac{(e^{u})^{y}}{1+e^{u}}
  \;=\;
  \tfrac12 e^{(y-\frac12)u}\!
  \int_0^{\infty}\!e^{-\frac12\omega u^{2}}\,PG(\omega\mid 1,0)\,d\omega,
  \qquad
  e^{-\frac12\omega c^{2}}PG(\omega\mid 1,0)\;\propto\;PG(\omega\mid 1,c).
\]
Conditionals then become
\[
  \omega_i\mid f,y \;\sim\; PG(1,f(x_i)),
  \qquad
  f\mid\omega,y \;\sim\; \mathcal{N}\!\bigl(\Sigma\kappa,\,\Sigma\bigr),
\]
with
\[
  \Sigma \;=\; (K_\theta^{-1}+\Delta)^{-1},
  \qquad
  \Delta \;=\; \diag\!\bigl(\E_q[\omega_i]\bigr),
  \qquad
  \kappa_i \;=\; y_i - \tfrac12.
\]
We fit the model by mean-field variational inference with $q(f,\omega)=q(f)\prod_i q(\omega_i)$, $q(f)=\mathcal{N}(\mu,\Sigma)$, $q(\omega_i)=PG(1,c_i)$. Alternation gives rise to the E and M steps of Sections~\ref{sec:estep}--\ref{sec:mstep}.

\subsection{Notation and code-to-math dictionary}
\label{sec:notation}
The following symbols appear verbatim in \texttt{pg\_classifier.py}:
\begin{center}
\begin{tabular}{@{}l l l@{}}
\toprule
\textbf{Code} & \textbf{Math} & \textbf{Meaning}\\
\midrule
\texttt{spectral.ws}        & $D$                  & $\sqrt{h^{d}\hat k(\xi)}\in\R^{m}$\\
\texttt{spectral.ws2}       & $D^{2}$              & $h^{d}\hat k(\xi)\in\R^{m}$\\
\texttt{spectral.Dprime}    & $\partial_\theta D^{2}$ & shape $(m,p)$\\
\texttt{spectral.fwd}       & $F$                  & NUFFT type-2, features $\to$ data\\
\texttt{spectral.fadj}      & $F^{*}$              & NUFFT type-1, data $\to$ features\\
\texttt{variational.delta}  & $\Delta$             & $\diag\E_q[\omega]\in\R^{n}$\\
\texttt{variational.mean}   & $\mu$                & $=\Sigma\kappa\in\R^{n}$\\
\texttt{variational.sigma\_diag} & $\diag\Sigma$  & posterior variance vector\\
\texttt{kappa}              & $\kappa$             & effective observation, Bernoulli $y-\tfrac12$\\
\texttt{probes}             & $\{z^{(j)}\}$        & Rademacher probes in $\{\pm1\}^{n}$\\
\texttt{beta\_x}            & $\beta^{(\mu)}=F^{*}K^{-1}\mu$ & M-step Term I coefficients\\
\texttt{beta\_probes}       & $\beta^{(j)}=F^{*}K^{-1}\Sigma z^{(j)}$ & M-step Term II coefficients\\
\bottomrule
\end{tabular}
\end{center}
The only quantity that is pinned down by the data (and, for NB, by the shape $r$) is $\kappa$; everything else in this dictionary moves during variational optimization.

\section{E step}
\label{sec:estep}

Fix $\theta$ and update $(\mu,\diag\Sigma,\Delta)$.

\subsection{Mean-field updates}
Optimizing the ELBO over $q(f)$ with $q(\omega)$ fixed gives a Gaussian with natural parameters determined entirely by $\Delta$, namely
\[
  q(f) \;=\; \mathcal{N}(\mu,\Sigma),
  \qquad
  \Sigma \;=\; (K_\theta^{-1}+\Delta)^{-1},
  \qquad
  \mu \;=\; \Sigma\kappa.
\]
Optimizing over $q(\omega)$ with $q(f)$ fixed gives $q(\omega_i)=PG(1,c_i)$ with
\[
  c_i \;=\; \sqrt{\E_{q(f)}[f(x_i)^{2}]}
       \;=\; \sqrt{\mu_i^{2}+\Sigma_{ii}},
  \qquad
  \E_q[\omega_i] \;=\; \frac{b_i}{2c_i}\tanh\!\frac{c_i}{2},
\]
with $b_i=1$ (Bernoulli) or $b_i=y_i+r$ (NB).

Therefore each E-step sweep needs three quantities: $\mu$, $\diag\Sigma$, and (after the update) the new $\Delta$. The challenge is that $\Sigma$ is never formed: both $\mu$ and $\diag\Sigma$ must come from a matrix-free primitive that applies $\Sigma$ to vectors.

\subsection{The $\Sigma$-apply primitive}
\label{sec:sigma-apply}
We want a matrix-free routine for $v\mapsto\Sigma v$ that never forms $K$ or $\Sigma$. Using $K=FD^{2}F^{*}=\Phi\Phi^{*}$ with $\Phi=FD$, the following chain of equivalences reduces the $n\times n$ system $(K^{-1}+\Delta)v=z$ to an $m\times m$ Hermitian feature-space system:
\begin{align}
  (K^{-1}+\Delta)\,v \;=\; z
  &\;\Longleftrightarrow\;
  (I+K\Delta)\,v \;=\; K\,z
  \label{eq:chain-setup}\\
  &\;\Longleftrightarrow\;
  (I+\Phi\Phi^{*}\Delta)\,v \;=\; \Phi\Phi^{*}\,z
  \label{eq:chain-phi}\\
  &\;\Longleftrightarrow\;
  (I+\Phi\Phi^{*}\Delta)\,\Phi u \;=\; \Phi\Phi^{*}\,z
  \label{eq:chain-image}\\
  &\;\Longleftrightarrow\;
  \Phi\,\bigl[u + \Phi^{*}\Delta\Phi\,u\bigr] \;=\; \Phi\Phi^{*}\,z
  \label{eq:chain-factor}\\
  &\;\Longleftrightarrow\;
  (I+\Phi^{*}\Delta\Phi)\,u \;=\; \Phi^{*}z,
  \qquad v \;=\; \Phi u.
  \label{eq:chain-cancel}
\end{align}
The two nontrivial steps both rely on structure of $\Phi$:
\begin{itemize}[leftmargin=1.4em,itemsep=1pt]
  \item Step \eqref{eq:chain-image}: we restrict to $v=\Phi u$. This loses no solutions because $v\in\operatorname{range}(K)=\operatorname{range}(\Phi)$; rearranging \eqref{eq:chain-setup} as $v = K(z-\Delta v)$ shows directly that $v$ lies in the column space of $K$, which equals that of $\Phi$.
  \item Step \eqref{eq:chain-cancel}: we cancel the outer $\Phi$ on both sides. This is valid because $\Phi$ is injective on $\C^{m}$: the columns of $\Phi$ are the (weighted) evaluations $\{\sqrt{h^{d}\hat k(\xi_j)}\,e^{i x\cdot\xi_j}\}_{j=1}^{m}$ of linearly independent trig modes on the data, so $\Phi u = 0$ forces $u=0$.
\end{itemize}

The final feature-space system is
\[
  \boxed{\;
    (I_{m}+\Phi^{*}\Delta\Phi)\,u \;=\; \Phi^{*}z,
    \qquad
    \Sigma z \;=\; \Phi u.
  \;}
\]
Substituting $\Phi = FD$ and multiplying through by $D$ on both sides (equivalently, working with the whitened variable $u$ as-is) gives the Hermitian positive-definite form actually used in code:
\begin{equation}
  \underbrace{(I_{m} + D F^{*}\Delta F D)}_{A}\,u \;=\; D F^{*} z,
  \qquad
  \Sigma z \;=\; F\,D\,u.
  \label{eq:sigma-apply}
\end{equation}
$A$ is applied matrix-free in $\calO(n+m\log m)$ per iteration via two NUFFTs and an elementwise multiply by $\Delta$. CG converges in feature space using nothing but the operator $A$.

\begin{lstlisting}
def A_feat(u):                        # (I + D F* Delta F D) u
    return u + ws * fadj(delta * fwd(ws * u))

def sigma_apply(z):                   # z -> Sigma z
    rhs    = ws * fadj(z)             # D F* z
    u, _   = cg(A_feat, rhs, tol=cg_tol)
    return fwd(ws * u).real           # F D u
\end{lstlisting}

Both \texttt{fwd} and \texttt{fadj} are batched, so \texttt{sigma\_apply} accepts a stack of right-hand sides and solves them in one CG call.

\subsection{$F^{*}\Delta F$ is Toeplitz for any diagonal $\Delta$}
\label{sec:weighted-toeplitz}

The operator $B:=F^{*}\Delta F$ that sits inside the Hermitian system $A=I+DBD$ has Toeplitz structure \emph{regardless} of the weight vector $\Delta$. Because the Fourier modes are equispaced on a tensor grid, $\xi_k = h\,\bm{k}$ with $\bm{k}\in\{-(m-1)/2,\ldots,(m-1)/2\}^{d}$, the $(j,k)$ entry is
\[
  [F^{*}\Delta F]_{jk}
  \;=\;
  \sum_{i=1}^{n}\Delta_{ii}\,e^{i x_i\cdot(\xi_k-\xi_j)}
  \;=\;
  \sum_{i=1}^{n}\Delta_{ii}\,e^{i h\,x_i\cdot(\bm{k}-\bm{j})},
\]
which depends only on the lag $\bm{r}=\bm{k}-\bm{j}$, not on $\bm{j}$ and $\bm{k}$ separately. Constant along diagonals is the defining property of a Toeplitz matrix; in $d>1$ the result is a $d$-level block-Toeplitz operator.

Equivalently, the data points $(x_i,\Delta_{ii})$ get absorbed into the Toeplitz \emph{symbol}
\[
  \tau[\bm{r}]
  \;=\;
  \sum_{i=1}^{n}\Delta_{ii}\,e^{i h\,x_i\cdot\bm{r}},
  \qquad
  \bm{r}\in\{-(m-1),\ldots,m-1\}^{d},
\]
which is a vector of length $(2m-1)^{d}$. Assembling $\tau$ takes one Type-1 NUFFT of the weights $\Delta_{ii}$ on the stencil locations $\bm{r}$. Once $\tau$ is built, applying $B$ to any vector $u$ is an ordinary length-$(2m-1)^{d}$ FFT-based circular convolution at cost $\calO(m^{d}\log m)$ with \emph{no further NUFFTs}:
\[
  B u \;=\; \mathcal{F}^{-1}\!\bigl(\mathcal{F}(\tau)\odot\mathcal{F}(\tilde u)\bigr)\big|_{\text{first-}m\text{-entries}},
\]
where $\tilde u$ is $u$ zero-padded to length $(2m-1)^{d}$. This replaces the two NUFFTs inside each CG iteration with a single FFT convolution and is the code path enabled by the \texttt{use\_exact\_weighted\_toeplitz\_operator} flag in \texttt{\_make\_sigma\_apply} and \texttt{\_make\_feature\_space\_solver}.

\paragraph{Where this shows up.}
The same observation is what makes the hierarchical block structure of Section~\ref{sec:hier} tractable: each per-location operator $T_\ell=F_\ell^{*}\Delta_\ell F_\ell$ is Toeplitz by the same calculation, applied to the subset $(x_i,\delta_i)_{i\in I_\ell}$. And the Toeplitz preconditioner of Appendix~\ref{app:precond} is built by specializing the symbol $\tau$ to the scalar approximation $\Delta\approx\bar\delta I$, but one could also use the \emph{exact} symbol if desired---the Toeplitz structure is present for any $\Delta$.

\subsection{Hutchinson estimator for $\diag\Sigma$}
With independent Rademacher probes $z^{(j)}\in\{\pm1\}^{n}$,
\[
  \E\!\bigl[z^{(j)}\odot(\Sigma z^{(j)})\bigr] \;=\; \diag(\Sigma),
\]
so $\tfrac{1}{J}\sum_j z^{(j)}\odot(\Sigma z^{(j)})$ is unbiased for $\diag\Sigma$. By stacking probes together with $\kappa$ into one right-hand-side block
\[
  Z \;=\; [\,\kappa;\, z^{(1)};\,\ldots;\,z^{(J)}\,]\in\R^{(J+1)\times n},
\]
a single batched CG call on $A$ computes $\Sigma Z=[\mu;\,\Sigma z^{(1)};\ldots;\,\Sigma z^{(J)}]$ in one pass---the same CG iterations produce both the mean and the probe responses.

\subsection{PG update}
Once $\mu$ and $\diag\Sigma$ are in hand, compute $c_i^{2}=\mu_i^{2}+(\diag\Sigma)_i$ and
\[
  \Lambda_i \;=\; \frac{b_i}{2c_i}\tanh\!\frac{c_i}{2}.
\]
A damped natural-gradient update stabilizes CAVI:
\[
  \Delta^{(t+1)} \;=\; (1-\rho_t)\,\Delta^{(t)} + \rho_t\,\Lambda,
  \qquad
  \rho_t \;=\; \frac{\rho_0}{1+\gamma t},
  \qquad \Delta \;\ge\; 0.
\]

\begin{lstlisting}
Z          = cat([kappa[None, :], probes], dim=0)   # (1+J, n)
S_all      = sigma_apply(Z)                         # batched Sigma Z
mean       = S_all[0]                               # mu
Sz         = S_all[1:]                              # Sigma z^(j)
sigma_diag = (probes * Sz).mean(dim=0)              # Hutchinson

c2     = (sigma_diag + mean.pow(2)).clamp_min(1e-12)
c      = c2.sqrt()
Lambda = 0.5 * pg_b * torch.tanh(0.5 * c) / c       # PG posterior mean
rho    = rho0 / (1.0 + gamma * it)
delta.mul_(1.0 - rho).add_(rho * Lambda)            # damped update
delta.clamp_(min=0.0)
\end{lstlisting}

This is the entire inner E-step iteration. The only nontrivial work is one batched CG call on the Hermitian operator $A$ of \eqref{eq:sigma-apply}. A legacy Woodbury formulation is described in Appendix~\ref{app:woodbury}; it is superseded by the stable form above.

\section{M step}
\label{sec:mstep}

Hold $q(f)$ fixed and differentiate the ELBO in the kernel hyperparameters $\theta$. Only the prior term depends on $\theta$:
\[
  \mathcal{L}(\theta) \;=\;
  -\tfrac12\log|K_\theta|
  \;-\;\tfrac12\,\mu^{\top}K_\theta^{-1}\mu
  \;-\;\tfrac12\,\tr(K_\theta^{-1}\Sigma)
  \;+\;\text{const}.
\]

\subsection{From three terms to two}
\label{sec:revised-mstep}
Differentiating and combining traces gives
\[
  \partial_\theta\mathcal{L}
  \;=\;
  \tfrac12\,\mu^{\top}K^{-1}\partial_\theta K\,K^{-1}\mu
  \;+\;\tfrac12\,\tr\!\bigl(\Sigma\,K^{-1}\partial_\theta K\,K^{-1}\bigr)
  \;-\;\tfrac12\,\tr(K^{-1}\partial_\theta K).
\]
The two trace terms collapse. From $\Sigma(K^{-1}+\Delta)=I$ we get $\Sigma K^{-1}=I-\Sigma\Delta$, so
\[
  \tr\!\bigl(\Sigma K^{-1}\partial_\theta K\,K^{-1}\bigr)
  \;=\;
  \tr\!\bigl((I-\Sigma\Delta)\,\partial_\theta K\,K^{-1}\bigr)
  \;=\;
  \tr(K^{-1}\partial_\theta K)
  \;-\;
  \tr(K^{-1}\Sigma\Delta\,\partial_\theta K),
\]
and the $\tr(K^{-1}\partial_\theta K)$ term cancels against the $-\tfrac12\tr(K^{-1}\partial_\theta K)$ above. What remains is the \emph{revised two-term gradient}:
\begin{equation}
\boxed{\;
  \partial_\theta\mathcal{L}
  \;=\;
  \tfrac12\,\underbrace{\mu^{\top}K^{-1}\partial_\theta K\,K^{-1}\mu}_{\text{Term I}}
  \;-\;
  \tfrac12\,\underbrace{\tr\!\bigl(K^{-1}\Sigma\Delta\,\partial_\theta K\bigr)}_{\text{Term II}}.
\;}
\label{eq:mstep-grad}
\end{equation}
The original three-term form is derived in Appendix~\ref{app:three-term}; it is not used.

\subsection{Two linear systems}
\label{sec:mstep-two-systems}

The whole M step reduces to two linear-algebra primitives. Everything else is an NUFFT or a diagonal multiply, and $K^{-1}$ is never formed. Both primitives rewrite $K^{-1}v$ for a data-space vector $v$ using the same algebraic trick.

\paragraph{The trick.}
Because $\Sigma^{-1}=K^{-1}+\Delta$,
\begin{equation}
  K^{-1}v \;=\; \Sigma^{-1}v \;-\; \Delta v \qquad\text{for any }v.
  \label{eq:Kinv-trick}
\end{equation}
Because $K=FD^{2}F^{*}$, any $v\in\operatorname{range}(K)$ can be written as $v = FD^{2}\beta(v)$ with
\[
  \beta(v) \;:=\; F^{*}K^{-1}v,
\]
i.e.\ $\beta(v)$ is the Fourier spectrum that reconstructs $v$ through the spectral kernel model. Both primitives compute some $\beta(v)$ \emph{without} ever inverting $K$.

\subsubsection{System I: $K^{-1}\mu$ (for Term I)}

We want $\beta^{(\mu)}:=F^{*}K^{-1}\mu$. Plug $v=\mu$ into \eqref{eq:Kinv-trick} and use $\Sigma^{-1}\mu=\kappa$, which holds by construction because $\mu$ is defined as $\Sigma\kappa$:
\[
  K^{-1}\mu \;=\; \kappa - \Delta\mu.
\]
Apply $F^{*}$ and substitute $\mu=FD^{2}\beta^{(\mu)}$ to eliminate $\mu$:
\begin{equation}
  \beta^{(\mu)} \;=\; F^{*}\kappa \;-\; F^{*}\Delta\,F\,D^{2}\,\beta^{(\mu)}
  \quad\Longleftrightarrow\quad
  \boxed{\;(I_m + F^{*}\Delta F D^{2})\,\beta^{(\mu)} \;=\; F^{*}\kappa.\;}
  \label{eq:sysI}
\end{equation}
One deterministic linear solve for one vector. No probes, no stochasticity.

Given $\beta^{(\mu)}$, Term I follows directly. Using $\partial_\theta K = F(\partial_\theta D^{2})F^{*}$ and $\beta^{(\mu)}=F^{*}K^{-1}\mu$,
\begin{equation}
  \mu^{\top}K^{-1}\partial_\theta K\,K^{-1}\mu
  \;=\;
  \bigl(\beta^{(\mu)}\bigr)^{*}(\partial_\theta D^{2})\,\beta^{(\mu)}
  \;=\;
  \sum_{k=1}^{m}(\partial_\theta D^{2})_{k}\,\bigl|\beta^{(\mu)}_{k}\bigr|^{2}.
  \label{eq:termI}
\end{equation}

\subsubsection{System II: $K^{-1}\Sigma z^{(j)}$ (for Term II)}

Term II is a trace; we estimate it with Rademacher probes $z^{(j)}\in\{\pm1\}^{n}$,
\[
  \tr\!\bigl(K^{-1}\Sigma\Delta\,\partial_\theta K\bigr)
  \;=\;
  \tr\!\bigl(\Delta\,\partial_\theta K\,K^{-1}\Sigma\bigr)
  \;=\;
  \E_z\!\bigl[(\Delta z)^{\top}\partial_\theta K\,K^{-1}\Sigma\,z\bigr].
\]
Evaluating the integrand requires $F^{*}K^{-1}\Sigma z^{(j)}$ for each probe. Apply the trick \eqref{eq:Kinv-trick} with $v=\Sigma z^{(j)}$ and use $\Sigma^{-1}\Sigma z^{(j)}=z^{(j)}$:
\[
  K^{-1}\Sigma z^{(j)} \;=\; z^{(j)} - \Delta\Sigma z^{(j)}.
\]
Apply $F^{*}$ and substitute $\Sigma z^{(j)}=FD^{2}\beta^{(j)}$:
\begin{equation}
  \boxed{\;(I_m + F^{*}\Delta F D^{2})\,\beta^{(j)} \;=\; F^{*}z^{(j)},
    \qquad j=1,\ldots,J.\;}
  \label{eq:sysII}
\end{equation}
$J$ linear solves, one per probe.

Substituting $\partial_\theta K=F(\partial_\theta D^{2})F^{*}$ and defining $R^{(j)}:=F^{*}(\Delta z^{(j)})$,
\[
  (\Delta z^{(j)})^{\top}\partial_\theta K\,K^{-1}\Sigma z^{(j)}
  \;=\;
  (R^{(j)})^{*}(\partial_\theta D^{2})\,\beta^{(j)}
  \;=\;
  \sum_{k}(\partial_\theta D^{2})_{k}\,\overline{R^{(j)}_{k}}\,\beta^{(j)}_{k}.
\]
Averaging over probes and taking the real part (the trace is real) gives the Hutchinson estimate of Term II:
\begin{equation}
  \tr\!\bigl(K^{-1}\Sigma\Delta\,\partial_\theta K\bigr)
  \;\approx\;
  \frac{1}{J}\sum_{j=1}^{J}\,\Re\!\sum_{k=1}^{m}
  (\partial_\theta D^{2})_{k}\,\overline{R^{(j)}_{k}}\,\beta^{(j)}_{k}.
  \label{eq:termII}
\end{equation}

\subsubsection{Shared Hermitian operator and batched CG}

Systems \eqref{eq:sysI} and \eqref{eq:sysII} have the \emph{same} left-hand side but different right-hand sides, so one CG call solves them all together. The operator $I+F^{*}\Delta F D^{2}$ is not Hermitian (the diagonal $D^{2}$ sits only on one side), which is bad for CG, so we symmetrize with the substitution $y:=D\beta$. Multiplying \eqref{eq:sysI}--\eqref{eq:sysII} by $D$ on the left gives the Hermitian positive definite system
\begin{equation}
  A\,y \;=\; D\,q,
  \qquad
  A \;:=\; I_m + D F^{*}\Delta F D,
  \qquad
  \beta \;=\; y / D,
  \label{eq:Asys}
\end{equation}
with $q=F^{*}\kappa$ for System I and $q=F^{*}z^{(j)}$ for System II. The operator $A$ is \emph{identical} to the $\Sigma$-apply operator of \eqref{eq:sigma-apply}, so the same CG machinery is reused.

Plugging \eqref{eq:termI} and \eqref{eq:termII} into the gradient \eqref{eq:mstep-grad} gives, for each hyperparameter $\ell=1,\ldots,p$,
\begin{equation}
\boxed{\;
  \bigl[\partial_\theta\mathcal{L}\bigr]_\ell
  \;=\;
  \tfrac12\sum_{k=1}^{m}(\partial_\theta D^{2})_{k\ell}\,\bigl|\beta^{(\mu)}_{k}\bigr|^{2}
  \;-\;
  \frac{1}{2J}\sum_{j=1}^{J}\Re\!\sum_{k=1}^{m}
  (\partial_\theta D^{2})_{k\ell}\,\overline{R^{(j)}_{k}}\,\beta^{(j)}_{k}.
\;}
\label{eq:mstep-final}
\end{equation}

\subsection{Implementation}
\label{sec:mstep-impl}

Inputs: $\kappa\in\R^{n}$, $\Delta\in\R^{n}$, $D\in\R^{m}$ (\texttt{ws}), $\partial_\theta D^{2}\in\R^{m\times p}$ (\texttt{Dprime}), operators $F=$\texttt{fwd}, $F^{*}=$\texttt{fadj}, probe count $J=$\texttt{n\_probes}.

\begin{enumerate}[leftmargin=1.6em,itemsep=2pt]
  \item \textbf{Draw probes.} Sample $z^{(1)},\ldots,z^{(J)}\in\{\pm1\}^{n}$.
  \item \textbf{Build all RHSs at once.} One batched $F^{*}$ applied to $[z^{(1)},\ldots,z^{(J)},\kappa]$ gives
    $Q=[F^{*}z^{(1)},\ldots,F^{*}z^{(J)},F^{*}\kappa]\in\C^{(J+1)\times m}$.
  \item \textbf{One batched CG.} Solve $A\,y_{\text{all}}=D Q$ for all $J{+}1$ right-hand sides simultaneously using $A u = u + D F^{*}(\Delta\,F(Du))$. Divide by $D$ to get $\beta_{\text{all}}$, and split into $\beta^{(j)}$ ($j=1,\ldots,J$) and $\beta^{(\mu)}$.
  \item \textbf{Term I.} Compute $|\beta^{(\mu)}|^{2}$ elementwise, then
    $\text{term1}=(\partial_\theta D^{2})^{\top}|\beta^{(\mu)}|^{2}\in\R^{p}$ (eq.~\eqref{eq:termI}).
  \item \textbf{Probe side of Term II.} One batched $F^{*}$ applied to $[\Delta z^{(1)},\ldots,\Delta z^{(J)}]$ gives $R\in\C^{J\times m}$.
  \item \textbf{Assemble Term II.} Form $X_{jk}=\overline{R^{(j)}_{k}}\,\beta^{(j)}_{k}$ and compute
    $\text{per\_probe}_{j\ell}=\Re\sum_{k}(\partial_\theta D^{2})_{k\ell}X_{jk}\in\R^{J\times p}$,
    then $\text{term2}=\tfrac{1}{J}\sum_j\text{per\_probe}_j\in\R^{p}$ (eq.~\eqref{eq:termII}).
  \item \textbf{Gradient.} $\text{grad}=\tfrac12(\text{term1}-\text{term2})\in\R^{p}$ (eq.~\eqref{eq:mstep-final}).
\end{enumerate}

\begin{lstlisting}
# Step 1: probes  z^(j) in {+/-1}^n
probes = rademacher((n_probes, n))

# Step 2: RHS block  Q = F* [probes; kappa]
Q_block = fadj_batched(probes)                       # (J, m)   F* z^(j)
q_y     = fadj_batched(kappa.unsqueeze(0))           # (1, m)   F* kappa
Q_all   = cat([Q_block, q_y], dim=0)                 # (J+1, m)

# Step 3: one batched CG on  A y = D Q,  A = I + D F* Delta F D
#         then beta = y / D.  solve_A_beta hides the symmetrization.
beta_all, _ = solve_A_beta(Q_all)                    # (J+1, m)
beta_probes = beta_all[:-1, :]                       # beta^(j),  (J, m)
beta_x      = beta_all[-1, :]                        # beta^(mu), (m,)

# Step 4: Term I  =  (dD2)^T |beta^(mu)|^2
abs2  = (beta_x.conj() * beta_x).real
term1 = Dprime.real.T @ abs2                         # (p,)

# Step 5: R^(j) = F* (Delta z^(j))
Rfeat = fadj_batched(apply_omega(probes.mT).T).T     # (J, m)

# Step 6: per-probe scalar  sum_k (dD2)_k conj(R^(j))_k beta^(j)_k, then mean
X     = Rfeat.conj() * beta_probes.T                 # elementwise, then sum
vals  = (X.mT @ Dprime).real                         # (J, p)
term2 = vals.mean(dim=0)                             # (p,)

# Step 7: gradient
grad = 0.5 * (term1 - term2)                         # (p,)
\end{lstlisting}

This is the entirety of \texttt{\_compute\_mstep\_gradient} in \texttt{pg\_classifier.py}. The dominant cost is one batched CG on $A$ with $J+1$ right-hand sides; everything else is a handful of NUFFTs and diagonal multiplies.

\section{Predictive mean}
\label{sec:pred-mean}

By the GP prior conditional, $p(f_{*}\mid f)=\mathcal{N}\!\bigl(k_{*}^{\top}K^{-1}f,\,k_{**}-k_{*}^{\top}K^{-1}k_{*}\bigr)$, so
\[
  p(f_{*}\mid y) \;=\; \int p(f_{*}\mid f)\,p(f\mid y)\,df
  \;\approx\; \int p(f_{*}\mid f)\,q(f)\,df,
\]
where the approximation replaces the intractable posterior $p(f\mid y)$ by its variational Gaussian $q(f)=\mathcal{N}(\mu,\Sigma)$. The integral is Gaussian in closed form, and taking its mean yields the latent predictive mean
\[
  \mu_{*} \;=\; \E_{q}[k_{*}^{\top}K^{-1}f]
  \;=\; k_{*}^{\top}K^{-1}\mu
  \;=\; F_{*}\,D^{2}\,(F^{*}K^{-1}\mu)
  \;=\; F_{*}\bigl(D^{2}\,\beta^{(\mu)}\bigr),
\]
where $F_{*}$ is the NUFFT type-2 built at $X_{*}$ and $\beta^{(\mu)}=F^{*}K^{-1}\mu$ is the same System I vector from the M step. One NUFFT type-2 on $D^{2}\beta^{(\mu)}$ gives the full vector of predictive means at all test points.

\begin{lstlisting}
pred_mean = eval_op.type2(ws2 * beta_mean)    # F_*(D^2 beta^(mu))
\end{lstlisting}

If the M step has just run, $\beta^{(\mu)}$ is already cached; otherwise it is obtained from one CG solve on the System I equation \eqref{eq:sysI}.

\section{Predictive variance}
\label{sec:pred-var}

The predictive variance comes from the \emph{law of total variance} applied to the same Gaussian integral $\int p(f_{*}\mid f)\,q(f)\,df$ of Section~\ref{sec:pred-mean}. For any Gaussian $q(f)=\mathcal{N}(\mu,\Sigma)$,
\begin{equation}
  \Var_{q}(f_{*})
  \;=\;
  \underbrace{\E_{q}\!\bigl[\Var(f_{*}\mid f)\bigr]}_{k_{**}-k_{*}^{\top}K^{-1}k_{*}}
  \;+\;
  \underbrace{\Var_{q}\!\bigl[\E(f_{*}\mid f)\bigr]}_{\Var_{q}[k_{*}^{\top}K^{-1}f]\,=\,k_{*}^{\top}K^{-1}\Sigma K^{-1}k_{*}},
  \label{eq:total-var}
\end{equation}
because $p(f_{*}\mid f)=\mathcal{N}(k_{*}^{\top}K^{-1}f,\,k_{**}-k_{*}^{\top}K^{-1}k_{*})$ has an $f$-independent conditional variance. This is the same identity that yields the conjugate-case predictive variance: there one plugs in $\Sigma=K-K(K+\sigma^{2}I)^{-1}K$ (the true Gaussian-likelihood posterior covariance) and everything collapses to $k_{**}-k_{*}^{\top}(K+\sigma^{2}I)^{-1}k_{*}$. Here we plug in $\Sigma=(K^{-1}+\Delta)^{-1}$ instead, giving
\[
  \sigma_{*}^{2}
  \;=\;
  k_{**} \;-\; k_{*}^{\top}K^{-1}k_{*} \;+\; k_{*}^{\top}K^{-1}\Sigma K^{-1}k_{*}.
\]
Using $K^{-1}\Sigma K^{-1}=K^{-1}-(I+\Delta K)^{-1}\Delta$ (the same algebraic identity that drove the revised M step), the $K^{-1}$ terms cancel and
\[
  \boxed{\;
    \sigma_{*}^{2}
    \;=\;
    k_{**} \;-\; k_{*}^{\top}(I+\Delta K)^{-1}\Delta\,k_{*}.
  \;}
\]
Equivalently, in the feature notation $\psi_{*}=D\varphi_{*}$ with $\varphi_{*}$ the unweighted Fourier row and $k_{*}=\Phi\psi_{*}$, the predictive variance becomes
\[
  \sigma_{*}^{2} \;=\; \psi_{*}^{*}(I+B)^{-1}\psi_{*},
  \qquad B \;=\; D F^{*}\Delta F D.
\]
The operator $I+B$ is $A$ from \eqref{eq:Asys}: the E step, M step, and predictive variance all reuse the same Hermitian feature-space solver.

\subsection{Exact per-test-point variance}
For each test point, solve
\[
  A\,u_{*} \;=\; \psi_{*},
  \qquad
  \sigma_{*}^{2} \;=\; \psi_{*}^{*}u_{*}.
\]
The code batches these solves across test points when possible; otherwise it loops (routine \texttt{\_predictive\_latent\_moments}). Exact but $\calO(n_{\text{test}})$ CG calls.

\subsection{Stochastic predictive variance via Rademacher probes}
\label{sec:pred-var-stoch}
For many test points this is too expensive. The trick: write the quadratic form as a sum over offset diagonals of $G:=DA^{-1}D$,
\[
  \sigma_{*}^{2}
  \;=\;
  \varphi_{*}^{*}G\,\varphi_{*}
  \;=\;
  \sum_{r\in\{-(m{-}1),\ldots,m{-}1\}^{d}}
  e^{2\pi i h\langle x_{*},r\rangle}\,c[r],
  \qquad
  c[r] \;:=\; \sum_{m_{\ell}-m_{j}=r}G_{j\ell}.
\]
The offsets $c[r]$ are estimated by Hutchinson probes $\eta^{(s)}\in\{\pm1\}^{m}$ in \emph{feature space}: solve $A\beta^{(s)}=D\eta^{(s)}$ via batched CG, form $\gamma^{(s)}=D\beta^{(s)}=G\eta^{(s)}$, and compute
\[
  \widehat c^{(s)}[r] \;=\; \sum_{j}\gamma^{(s)}_{j}\eta^{(s)}_{j+r},
\]
an unbiased estimator because $\E[\eta^{(s)}_{k}\eta^{(s)}_{\ell}]=\delta_{k\ell}$. For each probe this is a cross-correlation, evaluated by one FFT after zero-padding to shape $(2m-1)^{d}$. Averaging $J$ probes and then evaluating
\[
  \sigma_{*}^{2}\;\approx\;\sum_{r}e^{2\pi i h\langle x_{*},r\rangle}\,\widehat c[r]
\]
at all test points takes one final NUFFT type-2 on the $(2m-1)^{d}$ offset grid. The dominant cost is $J$ probe solves on $A$, independent of $n_{\text{test}}$.

\section{Predictive class probability}
\label{sec:pred-prob}
The Bernoulli response is
\[
  p(y_{*}=1\mid y)
  \;=\;
  \int \sigma(f_{*})\,\mathcal{N}(f_{*}\mid\mu_{*},\sigma_{*}^{2})\,df_{*}.
\]
We use the probit-style approximation
\[
  p(y_{*}=1\mid y)
  \;\approx\;
  \sigma\!\left(\frac{\mu_{*}}{\sqrt{1+\pi\sigma_{*}^{2}/8}}\right),
\]
implemented in \texttt{approximate\_logistic\_gaussian\_prob}.

\section{Negative binomial likelihood}
\label{sec:nb}

For negative binomial observations $y_i\mid f_i,r\sim\mathrm{NB}(r,\sigma(f_i))$, the PG augmentation applies with auxiliary variables $\omega_i\sim PG(y_i+r,0)$ and effective observation
\[
  \kappa_i \;=\; \frac{y_i-r}{2},
  \qquad
  b_i \;=\; y_i + r,
  \qquad
  \Lambda_{ii} \;=\; \frac{b_i}{2c_i}\tanh\!\frac{c_i}{2}.
\]
Every piece of the Fourier-feature linear algebra from Sections~\ref{sec:estep}--\ref{sec:pred-var} is unchanged; only the likelihood-level scalars $\kappa$, $b$, and $\Lambda$ differ from the Bernoulli case.

\subsection{Learning the shape parameter $r$}

When $r$ is also optimized, the only $r$-dependent part of the ELBO is the expected log-likelihood; differentiating it gives
\begin{equation}
\boxed{\;
  \frac{\partial\mathcal{L}}{\partial r}
  \;=\;
  \sum_{i=1}^{n}\Bigl[
    \psi(y_i+r) - \psi(r) + \E_{q(f_i)}\log\sigma(-f_i)
  \Bigr],
\;}
\end{equation}
where $\psi$ is the digamma function. The expectation $\E_{q(f_i)}\log\sigma(-f_i)$ is a one-dimensional Gaussian integral against $\mathcal{N}(\mu_i,\Sigma_{ii})$ and is evaluated by Gauss--Hermite quadrature. With standard nodes/weights $(\tilde t_k,\tilde w_k)$ normalized so that $\sum_k\tilde w_k h(\tilde t_k)\approx\int h(t)\,\mathcal{N}(t\mid 0,1)\,dt$,
\[
  \E_{q(f_i)}\log\sigma(-f_i)
  \;\approx\;
  \sum_{k=1}^{K}\tilde w_k\,
  \log\sigma\!\bigl(-(\mu_i+\sqrt{\Sigma_{ii}}\,\tilde t_k)\bigr).
\]
The whole quadrature step is vectorized as an $n\times K$ broadcast and costs $\calO(nK)$, with no NUFFT or CG solves once $\mu_i$ and $\Sigma_{ii}$ are in hand. To enforce positivity we optimize $\eta=\log r$:
\[
  \eta^{(t+1)} \;=\; \eta^{(t)} + \rho_r\,r^{(t)}\,\partial_r\mathcal{L}.
\]
After updating $r$, the next E step refreshes $\kappa$, $b$, and $\Lambda$; the expensive linear algebra is unchanged.

\subsection{Predictive mean count}
Under the mean-value parameterization $\E[y_i\mid f_i,r]=re^{f_i}$,
\[
  \E[y_{*}\mid y] \;\approx\; r\,e^{\mu_{*}+\sigma_{*}^{2}/2},
\]
which reuses exactly the latent predictive mean and variance from Sections~\ref{sec:pred-mean}--\ref{sec:pred-var}.

\section{Hierarchical additive kernel}
\label{sec:hier}

The hierarchical model (implemented separately in \texttt{polyagamma\_classification/hierarchical/}) extends the single-kernel setup to $L$ locations indexed by $\ell_i\in\{1,\ldots,L\}$, with
\[
  f_i \;=\; g(t_i) + h_{\ell_i}(t_i),
  \qquad g\sim GP(0,k_g),
  \quad h_\ell\stackrel{\text{ind}}{\sim}GP(0,k_h),
\]
so the prior covariance splits as $K=K_g+K_h$ with $K_g$ dense over all $n$ points and $K_h$ block-diagonal by location.

The PG variational updates carry over unchanged with the replacement $K_\theta\to K_g+K_h$. The only challenge is applying the block kernel efficiently. Both kernels admit equispaced Fourier feature representations on a shared spectral grid (take whichever is finer), giving
\[
  K_g = \Phi_g\Phi_g^{*},
  \qquad
  K_h = \Phi_{\text{loc}}\Phi_{\text{loc}}^{*},
  \qquad
  \Phi_{\text{loc}} \;=\; \operatorname{blockdiag}(\Phi_1,\ldots,\Phi_L).
\]
Stack them into $U=[\Phi_g\;\;\Phi_{\text{loc}}]\in\C^{n\times(1+L)m}$ so that $K=UU^{*}$. The feature-space reduction from Section~\ref{sec:sigma-apply} generalizes verbatim:
\[
  \boxed{\;\Sigma z \;=\; U(I+U^{*}\Delta U)^{-1}U^{*}z.\;}
\]
The reduced operator $I+U^{*}\Delta U$ has dimension $(1+L)m$ with block structure
\[
  I + U^{*}\Delta U
  \;=\;
  \begin{bmatrix}
    I + \Phi_g^{*}\Delta\Phi_g & \Phi_g^{*}\Delta\Phi_{\text{loc}}\\[2pt]
    \Phi_{\text{loc}}^{*}\Delta\Phi_g & I + \Phi_{\text{loc}}^{*}\Delta\Phi_{\text{loc}}
  \end{bmatrix}.
\]

\subsection{Toeplitz blocks and block CG matvec}

Let $I_\ell$ denote the index set of location $\ell$ and $\Delta_\ell=\diag(\delta_i:i\in I_\ell)$. The per-location weighted operator $T_\ell:=F_\ell^{*}\Delta_\ell F_\ell$ is Toeplitz by exactly the observation of Section~\ref{sec:weighted-toeplitz}: the symbol $\tau_\ell[\bm{r}]=\sum_{i\in I_\ell}\delta_i\,e^{i h x_i\cdot\bm{r}}$ is built by one Type-1 NUFFT on the subset $(x_i,\delta_i)_{i\in I_\ell}$, and $T_\ell$ is then applied as an FFT convolution. The global version $T_{\text{all}}:=F_g^{*}\Delta F_g=\sum_\ell T_\ell$ is the same construction on the full data. Partition the reduced-space vector $a=(a_g,a_1,\ldots,a_L)$ with $a_g\in\C^{m}$ and $a_\ell\in\C^{m}$; then one matvec with $I+U^{*}\Delta U$ produces
\[
  \bigl(I+U^{*}\Delta U\bigr)_g a
  \;=\;
  a_g + D_g T_{\text{all}}(D_g a_g) + \sum_{\ell=1}^{L}D_g T_\ell(D_h a_\ell),
\]
\[
  \bigl(I+U^{*}\Delta U\bigr)_\ell a
  \;=\;
  a_\ell + D_h T_\ell(D_h a_\ell) + D_h T_\ell(D_g a_g).
\]
Each CG iteration is one global Toeplitz apply plus $L$ local Toeplitz applies, for total cost $\calO\!\bigl((1+L)m\log m\bigr)$.

\subsection{Hermitian symmetrization and block-diagonal preconditioner}

As in the flat case we whiten with $D_g^{1/2}$, $D_h^{1/2}$ substitutions so that the diagonal blocks become Hermitian positive definite:
\[
  I + D_g^{1/2}T_{\text{all}}D_g^{1/2},
  \qquad
  I + D_h^{1/2}T_\ell D_h^{1/2},
\]
matching the single-kernel $I+DF^{*}\Delta FD$ form. Dropping the off-diagonal cross blocks yields a block-diagonal preconditioner whose application is $1+L$ independent Toeplitz CG solves, each cheap because it involves no NUFFTs.

\subsection{Posterior decomposition}
After convergence, $K^{-1}\mu=\kappa-\Delta\mu$ from the same algebraic identity used in Section~\ref{sec:mstep-two-systems}, so the posterior decomposes into global and local pieces for free:
\[
  \mu_g \;=\; \Phi_g\Phi_g^{*}(\kappa-\Delta\mu),
  \qquad
  \mu_{h,\ell} \;=\; \Phi_\ell\Phi_\ell^{*}(\kappa-\Delta\mu)_{I_\ell}.
\]
No additional solves are needed---only NUFFT forward and adjoint applies.

The complete hierarchical pipeline (building $T_\ell$, running block CG, the block-diagonal preconditioner, and the block Schur complement) is implemented in \texttt{pg\_hierarchical.py} and \texttt{pg\_hierarchical\_blockcd.py}; refer to those for the full code paths.

\appendix

\section{Deprecated and alternate formulations}
\label{app:deprecated}

This appendix collects derivations that were used in earlier drafts or experimental code paths but are \emph{not} in the current \texttt{pg\_classifier.py} main loop. They are retained for reference.

\subsection{Original three-term M-step gradient}
\label{app:three-term}
Before the trace collapse of Section~\ref{sec:revised-mstep}, the gradient reads
\[
  \partial_\theta\mathcal{L}
  \;=\;
  \underbrace{\tfrac12\,\mu^{\top}K^{-1}\partial_\theta K\,K^{-1}\mu}_{\text{I}}
  \;+\;
  \underbrace{\tfrac12\,\tr\!\bigl(\Sigma K^{-1}\partial_\theta K\,K^{-1}\bigr)}_{\text{II}}
  \;-\;
  \underbrace{\tfrac12\,\tr(K^{-1}\partial_\theta K)}_{\text{III}}.
\]
Term I is identical to what the current code uses. Term II was previously handled via a separate Woodbury identity
\[
  K^{-1}\Sigma K^{-1} \;=\; K^{-1} - (I+\Delta K)^{-1}\Delta,
\]
which then required estimating $\Tr\!\bigl(K^{-1}\partial_\theta K\bigr)$ by independent Hutchinson probes. Collapsing Terms II and III via $\Sigma K^{-1}=I-\Sigma\Delta$ (Section~\ref{sec:revised-mstep}) avoids that second trace estimator and cancels the log-determinant term analytically, which is both cheaper and more stable.

\subsection{CAVI warm-up and natural-gradient updates}
\label{app:cavi}

The derivation in Section~\ref{sec:estep} takes the fixed-point E-step form directly. An earlier draft framed the same updates as natural-gradient updates on the exponential-family parameters $\eta_1=\Sigma^{-1}\mu$, $\eta_2=-\tfrac12\Sigma^{-1}$ with minibatches of size $s$,
\[
  \eta_1^{(t)} \;\leftarrow\;
  (1-\rho_t)\eta_1^{(t-1)} + \rho_t\tfrac{n}{s}\,I_S^{\top}(y_S-\tfrac12),
\]
\[
  \eta_2^{(t)} \;\leftarrow\;
  (1-\rho_t)\eta_2^{(t-1)} - \tfrac{\rho_t}{2}\bigl[K_\theta^{-1}+\tfrac{n}{s}\,I_S^{\top}\Lambda_S I_S\bigr].
\]
In the non-sparse (no inducing points) case, $\eta_2$ stores a full $n\times n$ matrix, which is impractical. Reformulating the update in terms of $\Delta$ alone,
\[
  \Delta^{(t)} \;\leftarrow\; (1-\rho)\,\Delta^{(t-1)} + \rho\,\tfrac{n}{s}\,I_S^{\top}\Lambda_S I_S,
\]
keeps all the linear algebra in the Hermitian feature-space form of Section~\ref{sec:sigma-apply}. The current code uses the damped PG update of Section~\ref{sec:estep} with full batches and does not run a minibatch natural-gradient loop.

\subsection{Old Woodbury E-step with catastrophic cancellation}
\label{app:woodbury}

The earliest implementation expanded
\[
  \Sigma z
  \;=\;
  Kz \;-\; K(I+\Delta K)^{-1}\Delta K z,
\]
and solved in feature space via
\[
  q_E := \Phi^{*}\Delta K z,
  \qquad
  (I+\Phi^{*}\Delta\Phi)\,u_{\text{corr}} \;=\; q_E,
  \qquad
  \Sigma z \;=\; Kz - \Phi u_{\text{corr}}.
\]
The operator being inverted is the same Hermitian $A$ of \eqref{eq:sigma-apply}, but the right-hand side and reconstruction subtract two large vectors whose difference is the much smaller $\Sigma z$. At large $n$ with small $\Delta$, small CG solve error is amplified by catastrophic cancellation. The direct form
\[
  A u \;=\; D F^{*}z,
  \qquad
  \Sigma z \;=\; F D u
\]
in the main E step avoids the subtraction and is numerically stable.

\subsection{Toeplitz preconditioner and the scalar-$\Delta$ surrogate}
\label{app:precond}

Section~\ref{sec:weighted-toeplitz} shows that $F^{*}\Delta F$ is already Toeplitz for any diagonal $\Delta$, so the exact operator $A=I+DF^{*}\Delta FD$ can be applied as an FFT convolution. This is the \texttt{use\_exact\_weighted\_toeplitz\_operator} code path and it removes the NUFFTs from the inner CG iteration entirely.

An older variant instead constructs a cheaper \emph{approximate} Toeplitz surrogate by replacing $\Delta$ with its scalar average:
\[
  M_E \;:=\; I + \bar\delta\,D\,T_{0}\,D,
  \qquad
  \bar\delta \;:=\; \frac{1}{n}\sum_i\Delta_{ii},
  \qquad
  T_{0} \;:=\; F^{*}F.
\]
$T_{0}$ is the Toeplitz convolution operator with symbol $\tau_{0}[\bm{r}]=\sum_i e^{ih x_i\cdot\bm{r}}$, i.e.\ the same construction as in Section~\ref{sec:weighted-toeplitz} with all weights set to 1. $M_E$ can then be used either as a left preconditioner for CG on the exact $A$, or as the source of a warm-start initial iterate $x_0$ from $M_E x_0 = r_E$. Preconditioning is the strictly stronger version of warm-starting since it applies $M_E^{-1}$ at every outer iteration.

The surrogate is most effective when $\Delta$ is close to $\bar\delta I$; when $\Delta$ is highly non-uniform, using the \emph{exact} weighted Toeplitz symbol $\tau[\bm{r}]=\sum_i\Delta_{ii}e^{ihx_i\cdot\bm{r}}$ is both cheaper per iteration and tighter. In practice the exact-Toeplitz path of Section~\ref{sec:weighted-toeplitz} has subsumed this approximate preconditioner mode.

\subsection{Cheaper Term II variant}
\label{app:cheap-term2}

Using $K^{-1}\Sigma\Delta = (I+\Delta K)^{-1}\Delta$ (same algebra as in the trace collapse), Term II can be rewritten so that the probes never need to pass through $K^{-1}\Sigma$. The current code instead reuses the feature-space solves it already performs for System I and pays for the extra $F^{*}$, because the numerical stability of the $\beta^{(j)}$ form is better.

\end{document}
